% from ega3-1-fr.tex
\oldpage[III]{126}

est un isomorphisme. Enfin, la filtration sur $H^n(X,\mathcal F)$, définie
par les noyaux des homomorphismes canoniques
$H^n(X,\mathcal F)\to H^n(X,\mathcal F_k)$, est $\mathfrak J$-bonne
(0, 13.7.7) et $\varphi_n$ est un isomorphisme topologique.
\footnote{La démonstration qui suit, plus simple que la démonstration
initiale, et le complément relatif à la filtration de
$H^n(X,\mathcal F)$, nous ont été communiqués par J.-P. Serre.}
\end{corollary}

\begin{proof}
L'entier $n\geq0$ étant fixé dans cette démonstration, nous poserons pour
simplifier
\[
\mathrm H=H^n(X,\mathcal F),\qquad
\mathrm H_k=H^n(X,\mathcal F_k)
\tag{4.1.7.3}\label{III.4.1.7.3-fr}
\]
et
\[
\mathrm R_k=\operatorname{Ker}(\mathrm H\longrightarrow\mathrm H_k),
\qquad\text{sous-$A$-module de $\mathrm H$.}
\tag{4.1.7.4}\label{III.4.1.7.4-fr}
\]
La suite exacte de cohomologie
\[
H^n(X,\mathfrak J^{k+1}\mathcal F)\longrightarrow H^n(X,\mathcal F)
\longrightarrow H^n(X,\mathcal F_k)\xrightarrow{\partial}
H^{n+1}(X,\mathfrak J^{k+1}\mathcal F)
\longrightarrow H^{n+1}(X,\mathcal F)
\]
montre que l'on a aussi
\[
\mathrm R_k=\operatorname{Im}\bigl(
H^n(X,\mathfrak J^{k+1}\mathcal F)\longrightarrow H^n(X,\mathcal F)
\bigr)~;
\]
nous poserons
\[
\begin{aligned}
\mathrm Q_k
&=\operatorname{Ker}\bigl(H^{n+1}(X,\mathfrak J^{k+1}\mathcal F)
\longrightarrow H^{n+1}(X,\mathcal F)\bigr)\\
&=\operatorname{Im}\bigl(H^n(X,\mathcal F_k)
\longrightarrow H^{n+1}(X,\mathfrak J^{k+1}\mathcal F)\bigr).
\end{aligned}
\tag{4.1.7.5}\label{III.4.1.7.5-fr}
\]
On a donc la suite exacte
\[
0\longrightarrow\mathrm R_k\longrightarrow\mathrm H
\longrightarrow\mathrm H_k\longrightarrow\mathrm Q_k\longrightarrow0.
\tag{4.1.7.6}\label{III.4.1.7.6-fr}
\]

\begin{env}[4.1.7.7]
\label{III.4.1.7.7-fr}
Soit $x$ un élément de $\mathfrak J^m$ ($m\geq0$)~; la multiplication par
$x$ dans $\mathfrak J^k\mathcal F$ est un homomorphisme
$\mathfrak J^k\mathcal F\to\mathfrak J^{k+m}\mathcal F$, et donne lieu
par suite à un homomorphisme
\[
\mu_{x,m}:H^n(X,\mathfrak J^k\mathcal F)
\longrightarrow H^n(X,\mathfrak J^{k+m}\mathcal F).
\tag{4.1.7.8}\label{III.4.1.7.8-fr}
\]
Si on désigne par $S$ la $A$-algèbre graduée
$\bigoplus_{k\geq0}\mathfrak J^k$, on sait que les multiplications
$\mu_{x,m}$ définissent sur
$\mathrm E=\bigoplus_{k\geq0}H^n(X,\mathfrak J^k\mathcal F)$ une
structure de module gradué de type fini sur l'anneau gradué $S$
(\hyperref[III.3.3.2-fr]{3.3.2}), qui est noethérien (II, 2.1.5).
\end{env}

\begin{lemma}[4.1.7.9]
\label{III.4.1.7.9-fr}
Les sous-modules $\mathrm R_k$ de $\mathrm H$ définissent sur $\mathrm H$
une filtration $\mathfrak J$-bonne.
\end{lemma}

\begin{proof}
En premier lieu, montrons que l'on a
\[
\mathfrak J^m\mathrm R_k\subset\mathrm R_{k+m},
\tag{4.1.7.10}\label{III.4.1.7.10-fr}
\]
la multiplication dans $\mathrm H=H^n(X,\mathcal F)$ par un élément
$x\in\mathfrak J^m$ étant donc l'application $\mu_{x,0}$. Pour tout
$x\in\mathfrak J^m$, le diagramme
\[
\xymatrix{
\mathfrak J^{k+1}\mathcal F\ar[r]^x\ar[d]&
\mathfrak J^{k+m+1}\mathcal F\ar[d]\\
\mathcal F\ar[r]_x&\mathcal F
}
\tag{4.1.7.11}\label{III.4.1.7.11-fr}
\]

